\documentclass[12pt]{article}
\usepackage[utf8]{inputenc}
\usepackage[T1]{fontenc}
\usepackage{hyperref}
\title{Prasuni Language Resources for K-12 Curriculum Development}
\author{Compiled by Claude Sonnet 5.5 (claude-sonnet-5-5) for Andr\'as Kornai}
\date{}
\begin{document}
\maketitle

\section{Introduction}
Prasuni (Prasun, Wasi-wari, Wasi-weri; ISO 639-3 \texttt{prn}) is a Nuristani language of the Indo-Iranian branch, spoken in the P\=ar\=un valley of northeastern Afghanistan by roughly 8,000 people, in three dialects across six villages. The Endangered Languages Project rates it vulnerable.

\section{Primary Educational Resources}
I found no primers, textbooks, or school programmes in Prasuni.

\section{Linguistic and Reference Materials}
\begin{itemize}
\item Georg Buddruss and Almuth Degener, \emph{Materialien zur Prasun-Sprache des afghanischen Hindukusch} (Harvard Oriental Series 80, 2016): texts with German translations and a glossary, and the basis for the first comprehensive grammar of Prasun. Listing: \url{https://media.standaardboekhandel.be/p/materialien-zur-prasun-sprache-des-afghanischen-hindukusch-9780674980372}
\item DFG project page for the Prasun materials: \url{https://gepris.dfg.de/gepris/projekt/271862156}
\item Morgenstierne's 1949 sketch, \emph{The Language of the Prasun Kafirs}, is cited in the search summary as the earlier main reference. I did not find an online copy.
\item Wikipedia entry on Prasuni: \url{https://Www.wikipedia.org/wiki/Prasuni}
\item Wiktionary Wasi-wari entry: \url{https://www.wiktionary.com/wiki/Wasi-wari}
\item WALS language page: \url{https://wals.info/languoid/lect/wals_code_pra}
\end{itemize}
I found no standalone dictionary.

\section{Religious and Community Resources}
None found in my searches.

\section{Digital and Audio Resources}
\begin{itemize}
\item Endangered Languages Project record: \url{https://endangeredlanguages.com/node/91336}
\item Database entry for the Prasun doculect: \url{https://lemontree210.pythonanywhere.com/en/doculect/prasun}
\end{itemize}
I did not verify any audio content.

\section{Curriculum Development Context}
The scholarly record is a grammar and texts for researchers, not school material. Any curriculum would need community-led work. The security and access situation in Afghanistan is a practical constraint that I did not research.

\end{document}